% ============================================================
% Convergence in Metric Spaces
% ============================================================
% =========================================================
\subsubsection{Convergence in metric space}
% =========================================================
\begin{definition}[Convergence in a Metric Space]
Let $(X,d)$ be a metric space and let $(x_n)$ be a sequence in $X$.
We say that $(x_n)$ \emph{converges} to a point $x \in X$ if

\[
\forall \varepsilon > 0 \;\exists N \in \mathbb{N}
\text{ such that }
n \ge N \Rightarrow d(x_n, x) < \varepsilon.
\]

In this case we write
\[
x_n \to x.
\]
\end{definition}

% ------------------------------------------------------------

\begin{lemma}[Convergence is Preserved by Isometries]
Let $(X,d_X)$ and $(Y,d_Y)$ be metric spaces, and let
$f : X \to Y$ be an isometry.

If $(x_n)$ is a sequence in $X$ and $x_n \to x$ in $X$,
then
\[
f(x_n) \to f(x)
\quad \text{in } Y.
\]
\end{lemma}



% ============================================================
% Closed Sets — Sequential Definition (Metric Version)
% ============================================================

\begin{definition}[Closed Set — Sequential Characterization]
Let $(X,d)$ be a metric space and let $F \subseteq X$.
The set $F$ is said to be \emph{closed} if whenever a sequence
$(x_n)$ in $F$ converges to a point $x \in X$, then $x \in F$.

Equivalently,
\[
\forall (x_n) \subseteq F,
\quad
x_n \to x \text{ in } X
\;\Longrightarrow\;
x \in F.
\]

In words: a closed set contains all limits of its convergent sequences.
\end{definition}



% ============================================================
% Structural Core of a Metric Space (Pre-Topology)
% ============================================================

\subsection*{Structural Core of a Metric Space}

Let $(X,d)$ be a metric space.

All analytic structure in $X$ is built from the metric $d$
via the following chain:

\[
\text{Metric}
\;\Longrightarrow\;
\text{Balls}
\;\Longrightarrow\;
\text{Convergence}
\;\Longrightarrow\;
\text{Closed Sets (Sequential Form)}.
\]

\paragraph{1. Metric $\Longrightarrow$ Balls}
Given $x \in X$ and $r>0$, define the open ball
\[
B_r(x) := \{ y \in X : d(x,y) < r \}.
\]
Thus balls are determined entirely by the distance function.

\paragraph{2. Balls $\Longrightarrow$ Convergence}
A sequence $(x_n)$ converges to $x \in X$ if
\[
\forall \varepsilon > 0,\;
\exists N \in \mathbb{N}
\text{ such that }
n \ge N \Rightarrow x_n \in B_\varepsilon(x).
\]
Thus convergence means eventual containment in every ball
centered at the limit.

\paragraph{3. Convergence $\Longrightarrow$ Closed Sets (Sequential Form)}
A subset $F \subseteq X$ is called closed if
whenever $(x_n) \subseteq F$ and $x_n \to x$ in $X$,
then $x \in F$.

Thus closed sets are exactly those sets that contain
all limits of convergent sequences they contain.

\bigskip

\noindent
\textbf{Summary.}
\begin{itemize}
\item The metric defines distance.
\item Distance defines balls.
\item Balls define convergence.
\item Convergence determines closedness (in metric spaces).
\end{itemize}

All analytic notions in a metric space flow from the
triangle inequality.


\begin{remark}[Balls, Convergence, and Closed Sets]
Let $(X,d)$ be a metric space.

\begin{itemize}
    \item \textbf{Balls measure convergence.}
    A sequence $(x_n)$ converges to $x$ if and only if
    for every $\varepsilon > 0$, all sufficiently large terms
    of the sequence lie inside the ball $B_\varepsilon(x)$.
    Thus convergence is entirely described in terms of
    eventual containment in arbitrarily small balls.

    \item \textbf{Closed sets contain limits.}
    A subset $F \subseteq X$ is closed (sequentially) if
    whenever $(x_n) \subseteq F$ and $x_n \to x$,
    then the limit $x$ also lies in $F$.
\end{itemize}

\noindent
In summary:
\[
\text{Balls describe how sequences approach a point;}
\quad
\text{closed sets describe where limits are allowed to remain.}
\]
\end{remark}

\begin{theorem}[Sequential Characterization of Closure Points]
Let $(X,d)$ be a metric space, let $A \subseteq X$,
and let $x \in X$.

There exists a sequence $(a_n)$ in $A$ such that
$a_n \to x$ if and only if
\[
\forall r > 0,\quad B_r(x) \cap A \neq \varnothing.
\]
\end{theorem}





\begin{theorem}[Closed Sets and Open Complements]
Let $(X,d)$ be a metric space and let $F \subseteq X$.
Then $F$ is closed if and only if its complement
\[
X \setminus F
\]
is open in $X$.
\end{theorem}


\begin{theorem}[Arbitrary Unions of Open Sets are Open]
Let $(X,d)$ be a metric space and let $\{U_i\}_{i \in I}$
be a collection of open subsets of $X$.
Then
\[
\bigcup_{i \in I} U_i
\]
is open in $X$.
\end{theorem}

\begin{theorem}[Finite Intersections of Open Sets are Open]
Let $(X,d)$ be a metric space and let $U_1,\dots,U_n$
be open subsets of $X$.
Then
\[
\bigcap_{k=1}^{n} U_k
\]
is open in $X$.
\end{theorem}



\begin{theorem}[Intersection of Closed Sets is Closed]
Let $(X,d)$ be a metric space and let $\{F_i\}_{i \in I}$ be a collection of
closed subsets of $X$. Then
\[
\bigcap_{i \in I} F_i
\]
is closed in $X$.
\end{theorem}

\begin{theorem}[Finite Union of Closed Sets is Closed]
Let $(X,d)$ be a metric space and let $F_1,\dots,F_n \subseteq X$ be closed.
Then
\[
\bigcup_{k=1}^n F_k
\]
is closed in $X$.
\end{theorem}


\begin{theorem}[Equivalent Characterizations of Closure]
Let $(X,d)$ be a metric space, let $A \subseteq X$, and let $x \in X$.
The following statements are equivalent:

\begin{enumerate}
    \item[(a)] $x \in \overline{A}$.
    \item[(b)] For every $r > 0$,
    \[
    B_r(x) \cap A \neq \varnothing.
    \]
    \item[(c)] There exists a sequence $(a_n)$ in $A$ such that
    \[
    a_n \to x.
    \]
\end{enumerate}
\end{theorem}


\begin{definition}[Limit Point]
Let $(X,d)$ be a metric space, let $A \subseteq X$, and let $x \in X$.
We say that $x$ is a \emph{limit point} of $A$ if
\[
x \in \overline{A \setminus \{x\}}.
\]
\end{definition}

\begin{remark}[Equivalent Characterization]
A point $x$ is a limit point of $A$ if and only if
\[
\forall r > 0,\quad
B_r(x) \cap (A \setminus \{x\}) \neq \varnothing.
\]
\end{remark}

\begin{theorem}[Continuity of the Metric]
Let $(X,d)$ be a metric space. 
If $(x_n)$ and $(y_n)$ are sequences in $X$ such that
\[
x_n \to x
\quad \text{and} \quad
y_n \to y,
\]
then
\[
d(x_n, y_n) \to d(x,y).
\]
\end{theorem}

\begin{theorem}[Continuity of the Metric]
Let $(X,d)$ be a metric space. 
If $(x_n)$ and $(y_n)$ are sequences in $X$ such that
\[
x_n \to x
\quad \text{and} \quad
y_n \to y,
\]
then
\[
d(x_n, y_n) \to d(x,y).
\]
\end{theorem}











